% Version and date. The measured numbers are in numbers.tex, which tools/numbers.py writes from
% the results files. A few numbers are typed here because no results file holds them; each
% says where it comes from. numbers.tex also sets \ifsecondrun: false while the results are
% those of the first run only.
\newcommand{\paperdate}{23 November 2026}
\newcommand{\paperversion}{1.5}
% Generated by tools/numbers.py from verify/results/ and refs/rows-counts.json.
% Do not edit by hand.
\newif\ifsecondrun
\secondrunfalse
\newcommand{\filecount}{175{,}187}% scan.txt: lean files
\newcommand{\linecount}{17{,}374{,}215}% scan.txt
\newcommand{\linemillions}{17}% computed
\newcommand{\archivebytes}{239{,}275{,}554}% scan.txt
\newcommand{\archivesha}{9631b4f020199e580ba86178cb56c6e9b09f690c2369a80dc4334aa93f49c7fc}% scan.txt
\newcommand{\archiveshaA}{9631b4f020199e580ba86178cb56c6e9}% first half, for display
\newcommand{\archiveshaB}{b09f690c2369a80dc4334aa93f49c7fc}% second half, for display
\newcommand{\otherfiles}{3}% scan.txt
\newcommand{\mainlines}{12}% scan.txt
\newcommand{\scanmacros}{87}% scan.txt
\newcommand{\scanelabs}{4}% scan.txt
\newcommand{\scanruncmd}{5}% scan.txt
\newcommand{\scanelabfiles}{3}% scan.txt: files listed under elab
\newcommand{\scanruncmdfiles}{5}% scan.txt: files listed under run_cmd
\newcommand{\scanpartial}{2}% scan.txt
\newcommand{\scandepth}{203}% scan.txt
\newcommand{\scanimports}{909{,}822}% scan.txt
\newcommand{\scanallmathlib}{136{,}003}% scan.txt
\newcommand{\buildhours}{13.4}% build-summary.txt: END minus START
\newcommand{\buildjobs}{183{,}893}% build-summary.txt
\newcommand{\buildthreads}{32}% build-summary.txt
\newcommand{\buildcpu}{AMD EPYC 9575F 64-Core Processor}% build-summary.txt
\newcommand{\leanversion}{4.33.1}% build-summary.txt
\newcommand{\mathlibrev}{0df444a360eaa60ab8c11dca51a86af692955474}% build-summary.txt
\newcommand{\conesize}{1{,}478{,}219}% cut-result.txt
\newcommand{\conefive}{1{,}478{,}243}% cut-result.txt
\newcommand{\conefiveextra}{24}% difference of the two
\newcommand{\conearchive}{1{,}265{,}256}% the lists: names of the archive in the cone
\newcommand{\conerest}{212{,}963}% CONE-SIZE minus the archive's names in the cone
\newcommand{\projnames}{1{,}905{,}922}% the lists, by name
\newcommand{\projnamesout}{640{,}666}% the lists, by name
\newcommand{\projtheorems}{1{,}548{,}656}% the lists, by name
\newcommand{\coneprojtheorems}{1{,}075{,}015}% the lists, by name
\newcommand{\unusedtheorems}{473{,}641}% difference
\newcommand{\repeatedtheorems}{39{,}597}% theorem declarations minus theorem names
\newcommand{\theoremshare}{69\%}% computed
\newcommand{\projdefs}{331{,}347}% the lists, by name
\newcommand{\coneprojdefs}{178{,}298}% the lists, by name
\newcommand{\defshare}{54\%}% computed
\newcommand{\projunsaferec}{621}% the lists: definitions whose name ends in _unsafe_rec
\newcommand{\projmodules}{175{,}014}% cut-result.txt
\newcommand{\coneprojmodules}{147{,}489}% cut-result.txt
\newcommand{\moduleshare}{84\%}% computed
\newcommand{\unusedmodules}{27{,}698}% Lean files that declare nothing of the cone
\newcommand{\emptymodules}{173}% Lean files that declare nothing
\newcommand{\hubtheorems}{198}% cut-result.txt
\newcommand{\projopaque}{2}% the lists: opaque constants of the archive
\newcommand{\coneprojopaque}{0}% the lists: of them in the cone
\newcommand{\projinductives}{8{,}509}% the lists: inductive types of the archive
\newcommand{\replaymodules}{186{,}200}% replay-summary.txt
\newcommand{\replayconstants}{2{,}701{,}132}% replay-summary.txt
\newcommand{\replaydecls}{2{,}695{,}292}% replay-summary.txt
\newcommand{\replayskipped}{5{,}840}% replay-summary.txt
\newcommand{\replaytheorems}{2{,}084{,}256}% replay-summary.txt
\newcommand{\replaydefs}{556{,}318}% replay-summary.txt
\newcommand{\replayopaque}{3{,}882}% replay-summary.txt
\newcommand{\replayaxioms}{7}% replay-summary.txt
\newcommand{\replayquot}{4}% computed: the constants replayed less those of the other kinds; checked to be 4
\newcommand{\replayinductives}{15{,}769}% replay-summary.txt
\newcommand{\replayblocks}{15{,}734}% replay-summary.txt
\newcommand{\replayctors}{19{,}159}% replay-summary.txt
\newcommand{\replayrecs}{15{,}897}% replay-summary.txt
\newcommand{\replayminutes}{57}% replay-summary.txt: seconds / 60
\newcommand{\replayrepeated}{40{,}595}% replay-summary.txt
\newcommand{\replayskippedrec}{5{,}008}% replay-skipped.txt
\newcommand{\replaymillions}{2.7}% computed
\newcommand{\replayprogramlines}{582}% verify/replay/ReplayMain.lean: lines, comments included
\newcommand{\firstprogramlines}{294}% results/replay-program-first-form.lean: lines
\newcommand{\hostileprogramsha}{b5e3596dde657bc8}% hostile.txt: first sixteen digits
\newcommand{\hostilemodules}{29}% hostile.txt
\newcommand{\hostilebad}{25}% hostile.txt: modules a replay must reject
\newcommand{\hostilecontrols}{4}% hostile.txt
\newcommand{\hostilecrafted}{11}% hostile/: modules with no source file
\newcommand{\hostilefreshrejects}{9}% hostile.txt
\newcommand{\hostiledefaultrejects}{7}% hostile.txt
\newcommand{\independentfreshrejects}{9}% independent-hostile.txt
\newcommand{\independentdefaultrejects}{7}% independent-hostile.txt
\newcommand{\independentfreshminutes}{77}% independent-fresh.txt
\newcommand{\stockfreshrzminutes}{11}% stock.txt
\newcommand{\stockfreshrzcpuminutes}{11}% stock.txt: processor time
\newcommand{\stockreplayrzdecls}{430{,}239}% stock.txt
\newcommand{\stockreplayrzminutes}{1}% stock.txt
\newcommand{\stockreplayrzcpuminutes}{8}% stock.txt: processor time
\newcommand{\stockfreshmathlibminutes}{36}% stock.txt
\newcommand{\stockfreshmathlibcpuminutes}{36}% stock.txt: processor time
\newcommand{\stockreplaymathlibdecls}{766{,}714}% stock.txt
\newcommand{\stockreplaymathlibminutes}{5}% stock.txt
\newcommand{\stockreplaymathlibcpuminutes}{31}% stock.txt: processor time
\newcommand{\stockexportedmodules}{10{,}498}% replay-Mathlib/replay-summary.txt
\newcommand{\stockexported}{644{,}284}% replay-Mathlib/replay-summary.txt
\newcommand{\stockpartialunmarked}{531}% replay-Mathlib/replay-summary.txt
\newcommand{\toolchainfiles}{17{,}493}% toolchain.txt
\newcommand{\literaturerows}{186}% refs/rows-counts.json
\newcommand{\rowsdistinct}{184}% refs/rows-counts.json
\newcommand{\rowsrepeated}{2}% computed
\newcommand{\rowspublished}{179}% computed
\newcommand{\rowsgeneral}{90}% refs/rows-counts.json
\newcommand{\rowsspecial}{86}% refs/rows-counts.json
\newcommand{\rowsother}{3}% refs/rows-counts.json
\newcommand{\rowsinternal}{3}% refs/rows-counts.json

% Billed by the provider for the machine of the first run, from its start to its end; not a
% measurement of the archive.
\newcommand{\firstruncost}{26.81}
% The first replay was started with LEAN_NUM_THREADS=28; this is our record of the command,
% and no results file of that run holds it.
\newcommand{\firstreplaythreads}{28}
% The cost of the second machine, \secondruncost, is computed in tools/numbers.py from the
% provider's hourly price and the machine's start, which are typed there, and the last time in
% file-times.txt.
% The memory limit of the first machine's container, in GB, as the provider listed it; no
% results file of that run holds it.
\newcommand{\firstmemorylimit}{128}

% The outcome of the kernel replay, in words (used in the section on verification).
\newcommand{\replaystatus}{%
The kernel accepted every declaration sent to it: \replaytheorems{} theorems, \replaydefs{}
definitions, \replayopaque{} opaque constants, \replayaxioms{} axioms, the quotient type with
its \replayquot{} constants, and \replayinductives{} inductive types in \replayblocks{} blocks.
\ifsecondrun
The comparison found every one of the \replaydecls{} constants of the files, among them
\replayctors{} constructors and \replayrecs{} recursors, equal to the kernel's, and nothing
else in the kernel's environment. There was no failure. The run took \replayminutes{} minutes
on \replaythreads{} threads.
\else
The \replayctors{} constructors and \replayrecs{} recursors in the files agree with the ones
the kernel generated. There was no failure. The run took \replayminutes{} minutes on
\firstreplaythreads{} threads.
\fi
The \replayaxioms{} axioms are all declared in Lean's core library; the archive declares none,
and the five theorems depend on three of them\ifsecondrun\else, by the recorded lists of
\cref{sec:axioms}\fi. In \replayrepeated{} cases a name is declared
again by a further module, which Lean's import accepts for theorems with the same statement;
every version was checked. The \replayskipped{} declarations not replayed are those that Lean
marks unsafe or partial: \replayskippedrec{} are the compiled forms of recursive definitions
and the rest are programs and the unsafe axioms that Lean's compiler uses. Lean's own replay
does not check them either, and outside Lean's
module system the kernel accepts no theorem or safe definition that refers to one;
\cref{sec:checker} says what holds inside it.}
